\section{Fan contacts along cell boundaries}
\label{sec:area_law_fan_side_contacts}

The labels assigned across cell sides in \cite[Section~11]{OpenAI2026AreaLaw}
are labels of elementary perimeter segments. The following equalities
identify precisely where the incident triangles can meet those boundaries.
Empty and single-point intersections are included.

% Source: 10-geometry.tex, prop:two-families, lines 299--323,
% revision adc7f1241b42e322a6451854ab7e4b4c146bf78a.
% These are boundary identities, not the full two-family proposition.

\begin{theorem}[A triangle meets its own sides through its base]
    \label{thm:area_law_fan_own_side_contact}
    \lean{TNLean.PEPS.AreaLaw.Geometry.cellFanPolygon_inter_dyadicCellSide}
    \leanok
    \uses{def:area_law_cell_fan,def:area_law_dyadic_cell_side}
    Fix an arbitrary translated dyadic square, a natural exponent,
    a signed cell index and an optional side-midpoint subdivision.
    For any actual triangle $P_i$ of the resulting fan, write
    $e_i=[a_i,b_i]$ for its elementary base. If $E$ is any whole side
    of its own square, then
    \begin{align}
        P_i\cap E=e_i\cap E.
    \end{align}
\end{theorem}
\begin{proof}\leanok
    \uses{def:area_law_cell_fan,thm:area_law_cell_fan_cover,
        thm:area_law_elementary_side_boundary_geometry}
    Let $Q$ be the half-open square and $c$ its center.
    The elementary base lies in $\overline Q$, and
    $P_i$ is the convex join of $\{c\}$ and $e_i$.
    Since $Q$ is convex and $c\in\operatorname{int}Q$, every point
    of $P_i$ outside $e_i$ lies in $\operatorname{int}Q$.
    A whole side has one coordinate fixed at a lower or upper endpoint
    of the square, and so is disjoint from its interior. This gives
    $P_i\cap E\subseteq e_i\cap E$; base containment gives the reverse inclusion.
\end{proof}

\begin{theorem}[Contact with a distinct actual cell passes through the base]
    \label{thm:area_law_fan_other_cell_contact}
    \lean{TNLean.PEPS.AreaLaw.Geometry.fineLayer_cellFanPolygon_inter_closedCell_eq}
    \leanok
    \uses{def:area_law_cell_fan,def:area_law_fine_layer_indices}
    Fix an arbitrary common origin, finite endpoint set and nonnegative
    integer neighborhood width. Let $z\in F_{k,\ell_k}$ and
    $w\in F_{h,\ell_h}$ be actual fine-cell indices with
    $(k,z)\ne(h,w)$. Choose an arbitrary optional midpoint mask
    for the cell $Q_{\ell_k}(z)$, and let $P_i$ and $e_i$ be any
    resulting triangle and its elementary base. Then
    \begin{align}
        P_i\cap\overline{Q_{\ell_h}(w)}
        =e_i\cap\overline{Q_{\ell_h}(w)}.
    \end{align}
    No lower scale bound or positive-length contact assumption is needed.
\end{theorem}
\begin{proof}\leanok
    \uses{thm:area_law_fine_cells_disjoint,thm:area_law_cell_fan_cover}
    The two actual half-open fine cells are disjoint. The open interior
    of the reference cell is therefore disjoint from the other cell
    and from its closure. As in the preceding proof, every point
    of $P_i$ outside its base lies in that interior. Such a point
    cannot belong to the other closed cell, proving one inclusion.
    The other follows from $e_i\subseteq P_i$.
\end{proof}

These exact identities preserve the locations of contacts when passing from
triangles to their bases. Interfaces between run regions in different cells
will be considered through shared nondegenerate segments, rather than through
two arbitrary common points.
